\scanpage{283}
\mspage{233}

à quoi en vérité \unsure{serait-il|servait-il} bon ?! Pour le prouver, factorisons \inserted{encore} $S\unsure{-|\to}S'$ en
\diagram{\begin{tikzcd}[ampersand replacement=\&,column sep=large]
 S \arrow[r,hook,"{\substack{\text{imm.}\\\text{ouv. de}\\\text{Dwyer}}}"'] \& S'_0 \arrow[r,"{\in W}"'] \& S',
\end{tikzcd}}
d'où le diagramme de 2 carrés cocartésiens
\diagram{\begin{tikzcd}[ampersand replacement=\&]
 S \arrow[r,"{Cof_W}"] \arrow[d,hook,"{Dw}"'] \& X \arrow[d,hook,"{Dw}"] \\
 S'_0 \arrow[r,"{Cof_W}"] \arrow[d,"{\in W}"'] \& X'_0 \arrow[d,"{\in W}"] \\
 S' \arrow[r,"{Cof_W}"] \& X',
\end{tikzcd}}
où $X'_0-X'$ est $\in W$ d'après l'hypothèse sur $X$ \unsure{sous|sur} $S$. Il faut donc vérifier :

a) Un carré cocartésien construit sur une flèche qui est imm. de Dwyer, est \deleted{\unsure{homotopiq.}} $W$-cocartésien

b) Un carré commutatif \deleted{\unsure{dans|deux}} duquel deux côtés opposés sont dans $W$, est $W$-cocartésien.

c) Le composé de deux carrés $W$-cocartésiens est $W$-cocartésien.

Admettons provisoirement b) et c), et prouvons a), pour un carré coc.
\diagram{\begin{tikzcd}[ampersand replacement=\&,column sep=small,row sep=small]
 U \arrow[r,hook,"i"] \arrow[d,"p"'] \& X \arrow[d,"q"] \\
 U' \arrow[r,hook,"{i'}"] \& X'
\end{tikzcd}}
, où $\unsure{i|i'}$ est imm. \inserted{ouv.} de Dwyer. On procède \unsure{encore|alors} comme \unsure{tantôt|d'hab.}, en factorisant $p$ en $U\overset{\unsure{j\in|j}}{\hookrightarrow}U'_0\xrightarrow{p'}U'$, où $j$ est imm. ouv. de Dwyer, \inserted{$p'\in W$} \deleted{\unsure{On}}. d'où un diagramme en deux carrés cocartésiens accolés

% archive sheet number: 283

